\chapter{Gastroparesis}

\section*{Definition and Overview}
Gastroparesis is a condition in which the stomach empties food too slowly, because the muscles and nerves that normally move food along are not working properly. The word means "paralysed stomach". It causes nausea, vomiting, early fullness, bloating, and abdominal pain after eating, and it can seriously affect nutrition and quality of life. The most common known cause is diabetes, especially long-standing or poorly controlled diabetes, which damages the nerves controlling the stomach. Other causes include stomach surgery, some medicines, and conditions such as Parkinson's disease, while in many cases no cause is found (idiopathic). There is no cure, but diet, medicines, and procedures control symptoms. This chapter explains gastroparesis, its causes, and its management.

\section*{Epidemiology}
The exact prevalence of gastroparesis is unknown, but it is estimated to affect around 1--2% of the population, with higher rates in people with diabetes: around a third to a half of people with long-standing type 1 or type 2 diabetes may have delayed gastric emptying, though many have no symptoms. It is more common in women, and most diagnosed cases are in people aged 30--50. Idiopathic gastroparesis is the most common category in people without diabetes, and it can follow viral infections. The condition causes significant illness, and its symptoms are frequently misattributed to reflux or functional dyspepsia.

\section*{Aetiology and Pathophysiology}
\begin{itemize}
    \item \textbf{Diabetic gastroparesis:} high blood sugar damages the vagus nerve and the stomach's own nerves, slowing emptying; poor control and long duration increase risk
    \item \textbf{Post-surgical gastroparesis:} after operations on the stomach and oesophagus, when the vagus nerve is affected
    \item \textbf{Idiopathic gastroparesis:} no cause found in many cases, sometimes following a viral infection
    \item \textbf{Other causes:} Parkinson's disease, scleroderma, thyroid disease, and some medicines, including opioids and some diabetes drugs
    \item \textbf{Mechanism:} the stomach's contractions are weak or uncoordinated, and the valve at its exit does not relax properly, so food stays in the stomach
    \item \textbf{Bezoars:} undigested food can clump into solid masses that obstruct the stomach
    \item \textbf{Effects:} nausea, vomiting, poor intake, weight loss, and variable blood sugar in diabetes
    \item \textbf{High blood sugar itself:} hyperglycaemia slows gastric emptying, creating a cycle in diabetes
\end{itemize}

\section*{Clinical Features}
\begin{itemize}
    \item Nausea, which is often chronic and severe
    \item Vomiting, sometimes of food eaten hours earlier
    \item Early fullness: feeling full after only a small meal
    \item Bloating and abdominal pain or discomfort
    \item Loss of appetite and unintentional weight loss
    \item Fluctuating blood sugar in diabetes, with unpredictable highs and lows
    \item Heartburn and regurgitation, which can coexist
    \item Symptoms that worsen after fatty meals and large meals
    \item In severe cases, dehydration and malnutrition
\end{itemize}

\section*{Differential Diagnosis}
The symptoms overlap with many conditions.
\begin{itemize}
    \item Functional dyspepsia, which is common and similar
    \item Gastro-oesophageal reflux disease
    \item Peptic ulcer disease
    \item Bowel obstruction, which is an emergency with severe pain and vomiting
    \item Diabetic complications and autonomic neuropathy
    \item Eating disorders, in the differential of vomiting and weight loss
    \item The emptying test confirms delayed emptying and distinguishes gastroparesis
\end{itemize}

\section*{When to Seek Medical Care}
Anyone with persistent nausea, vomiting, early fullness, or weight loss should be assessed by a doctor. People with diabetes and unpredictable blood sugars should discuss the possibility of gastroparesis.

\textbf{Call 000 (or local emergency services) for:}
\begin{itemize}
    \item Severe abdominal pain with vomiting and inability to pass wind or stool, which may indicate obstruction
    \item Vomiting blood or coffee-ground material
    \item Signs of severe dehydration: no urine, confusion, or collapse
    \item Very high or very low blood sugar with severe symptoms
    \item Any emergency symptoms
\end{itemize}

\section*{Diagnosis and Investigations}
\begin{itemize}
    \item \textbf{History and examination:} the symptom pattern, diabetes, medicines, and previous surgery
    \item \textbf{Gastric emptying study:} the key test, in which a radio-labelled meal is followed on a scanner over hours; delayed emptying confirms the diagnosis
    \item \textbf{Breath tests:} for gastric emptying as an alternative
    \item \textbf{Gastroscopy:} to exclude obstruction, ulcers, and other structural causes
    \item \textbf{Blood tests:} for diabetes control, thyroid disease, and nutritional status
    \item \textbf{Further tests:} oesophageal and stomach manometry in selected cases
\end{itemize}

\section*{Management}
\begin{itemize}
    \item \textbf{Dietary changes:} small, frequent meals; low-fat, low-fibre foods; pureed or liquid meals if needed; and eating slowly
    \item \textbf{Blood sugar control:} in diabetes, tighter glucose control improves emptying, with careful insulin adjustment because emptying is unpredictable
    \item \textbf{Medicines:} prokinetics such as metoclopramide and domperidone stimulate stomach contractions, and anti-nausea medicines
    \item \textbf{Review of medicines:} stop or reduce opioids and other drugs that slow emptying
    \item \textbf{Hydration and nutrition:} oral rehydration and, in severe cases, supplements or feeding tubes
    \item \textbf{Procedures:} pyloromyotomy (G-POEM), a keyhole procedure that widens the stomach outlet, in selected people
    \item \textbf{Gastric electrical stimulation:} an implanted device that reduces nausea in some people with severe disease
    \item \textbf{Bezoar management:} endoscopic removal when food masses form
    \item \textbf{Psychological support:} nausea and vomiting are distressing, and support helps coping
\end{itemize}

\section*{Complications}
\begin{itemize}
    \item Weight loss and malnutrition
    \item Dehydration and electrolyte disturbance
    \item Bezoars and, rarely, obstruction
    \item Severe fluctuation of blood sugar in diabetes, increasing the risk of diabetic complications
    \item Oesophagitis from repeated vomiting
    \item Reduced quality of life, work, and social participation
    \item Complications of treatments and procedures
\end{itemize}

\section*{Prognosis and Outcomes}
Gastroparesis is a chronic condition, but most people improve with dietary modification and medicines, and symptoms can be controlled for many people. Diabetic gastroparesis improves with better glucose control, and some cases, especially post-viral, improve spontaneously over time. Severe disease is harder to manage, but new procedures and devices help a substantial number of people. The condition does not usually shorten life, but it can be disabling, and outcomes are best with a multidisciplinary approach that includes diet, medicines, diabetes care, and psychology.

\section*{Prevention and Self-Care}
\begin{itemize}
    \item Manage diabetes carefully to protect the nerves that control the stomach
    \item Eat small, frequent, low-fat, low-fibre meals
    \item Chew food thoroughly and eat slowly
    \item Sit upright for an hour after meals
    \item Stay hydrated with frequent small sips
    \item Avoid alcohol, smoking, and opioids
    \item Take medicines as prescribed and attend review
    \item Seek urgent care for severe pain, vomiting, or dehydration
\end{itemize}

\section*{Sources and Further Reading}
The following reputable sources provide further information for healthcare students and patients seeking reliable, current advice about gastroparesis.
\begin{itemize}
    \item The Gut Foundation. \textit{Gastroparesis.} https://www.gutfoundation.org.au/
    \item Healthdirect Australia. \textit{Gastroparesis.} https://www.healthdirect.gov.au/
    \item Diabetes Australia. \textit{Gastroparesis and diabetes.} https://www.diabetesaustralia.com.au/
    \item International Foundation for Gastrointestinal Disorders. \textit{Gastroparesis.} https://www.iffgd.org/
    \item Mayo Clinic. \textit{Gastroparesis.} https://www.mayoclinic.org/
    \item MedlinePlus. \textit{Gastroparesis.} https://medlineplus.gov/
\end{itemize}
